% !TeX root = ..\studio-operation-guide.tex
\section{Call Park}

You can use call park to place a call on hold, and then retrieve the
call from another phone in the system (for example, a phone in another
office or conference room). You can park an active call by pressing the
call park key on the phone. If the call is parked successfully, the
response is either a voice prompt confirming that the call was parked,
or a visible prompt on the LCD screen.

\begin{quote}
	\textbf{Note}
	
	Call park is not available on all servers. Contact your system
	administrator for more information.
\end{quote}

\subsubsection{To use call park:}

\begin{enumerate}
	\item User on phone A places a call to phone B.
	\item User on phone A wants to take the call in a conference room for
		privacy, and so presses the call park key on phone A.
	\item User on phone A walks to an available conference room where the phone
		is designated as phone C. The user dials call park retrieve code to
		retrieve the parked call.

		The system establishes call between phone C and B.
\end{enumerate}

\begin{quote}
	\textbf{Note}
	
	The call park code and call park retrieve code are predefined on the
	system server. Contact your system administrator for more information.
	
	If the parked call is not retrieved within a period of time assigned by
	the system, the phone performing call park will receive a call back.
\end{quote}
